% !TeX root = ../../main.tex

Das Ziel der Experimente ist es die einzelnen, Retry systematisch zu testen und zu untersuchen, wie dieser sich unter
den verschiedenen Konfigurationen verhält.
Zusätzlich soll analysiert werden wie stark der Einfluss der verschiedenen Parameter ist.

Die Parameter sind die Anzahl an erlaubten Retries, \emph{Retry Count}, die Dauer des Wartens zwischen den einzelnen
Retries, \emph{Delay}, und die einzelnen Strategien.
Die Strategien, wie in Abschnitt~\ref{sec:retry} erwähnt, sind als Basis kein Retry, immediate Retry, Retry mit festem
Delay und der Retry mit exponential backoff + jitter Delay.

Bewertet werden soll das mit verschiedenen Kenngrößen.
Die wichtigste hierbei ist die Erfolgsrate, die sich berechnet, in dem alle erfolgreichen Transaktionen durch die
Gesamtanzahl an Transaktionen geteilt werden.
Eine weitere entscheidende Kenngröße ist die durchschnittliche Ausführungszeit der Transaktionen.
Die berechnet sich aus der Summe der \texttt{total\_ms} aller Transaktionen geteilt durch die Anzahl der Transaktionen.
Des Weiteren werden die Quartile berechnet.
Dabei wird der Fokus auf die 50\%, 95\%, 99\% und 99.9\% Quartile gelegt, um zu zeigen, wie viel Prozent der
Transaktionen bis zu welchem Zeitpunkt abgeschlossen haben.
Die \emph{retry\_distribution} zeigt, wie viele Transaktionen wie häufig einen Retry ausgeführt haben.
Sortiert wird bei den genannten Kenngrößen nach successful oder failed.

Eine weitere Kenngröße ist der \emph{retry\_overhead}, der sich durch die durchschnittliche Differenz zwischen der
Gesamtdauer \texttt{total\_ms} und dem besten Versuch \texttt{best\_ms} berechnet.
\[
	\mathrm{retry\_overhead} =
	\frac{1}{n}
	\sum_{i=1}^{n}
	\left(
		\mathrm{total}_{\mathrm{ms},i}
		-
		\mathrm{best}_{\mathrm{ms},i}
	\right),
	\quad n = \mathrm{total\_transactions}
\]

Dieser Wert ist besonders relevant, da er den zusätzlichen Zeitaufwand zeigt, der durch die Ausführung von Retries
entsteht.
Dadurch lässt sich bewerten, welchen Einfluss Wiederholungsversuche auf die Gesamtausführungszeit einer Transaktion
haben.


Die letzte Kenngröße ist mit, eine der wichtigsten, der \emph{throughput}.
Der Durchsatz zeigt an, wie viele erfolgreiche Transaktionen innerhalb einer bestimmten Zeit verarbeitet werden
können.
Dadurch ermöglicht er eine Bewertung der Leistungsfähigkeit des Systems.
Dieser berechnet sich aus der Anzahl erfolgreicher Versuche(N) geteilt durch die Dauer des Experiments.
\[
	\mathrm{throughput} =
	\frac{
		N_{\mathrm{success}}
	}{
		t_{\mathrm{finish}}-t_{\mathrm{start}}
	}
\]
Es wird durch die Dauer des Experiments geteilt, da diese fest definiert sind.
In den Experimenten sind die Parameter variable, aber wenn ein Experiment angefangen wurde, läuft das die Anzahl an
Iterationen durch.
Dadurch lässt sich der Durchsatz der einzelnen Experimente berechnen und durch weitere Sortierungsparameter wird der
Durchsatz für diesen Parameter berechnet.

Damit diese Experimente repräsentative Ergebnisse liefern, gibt es Rahmenbedingungen.
Die Fehlerszenarien sind fest definiert für alle Transaktionen.
Wenn Retry beispielsweise für Deadlocks getestet werden soll, dann erzeugen die Transaktionen auf jeden Fall auch
einen Deadlock.
Bei den anderen Fehlerklassen passiert genau das Gleiche, dadurch lässt sich das Verhalten unter reproduzierbaren
Bedingungen beobachten.

Die Kenngrößen sollen im Endeffekt zeigen, ob es Retry Strategien gibt, die bei den einzelnen Fehlerklassen bessere
Ergebnisse liefern als andere Strategien.
Zusätzlich soll anhand der Ergebnisse abgeleitet werden, wann Retry sinnvoll eingesetzt werden kann.